\par\medskip{\large\bfseries\color{mauPhieu} Bài 1. Không gian mẫu và biến cố}\par\smallskip
\textbf{1. Kiến thức cần nhớ}
\begin{itemize}
\item \textbf{Phép thử ngẫu nhiên} (gọi tắt là phép thử) là một thí nghiệm hay một hành động mà ta không đoán trước được kết quả, nhưng biết được tập hợp tất cả các kết quả có thể xảy ra.
\item \textbf{Không gian mẫu}, kí hiệu $\Omega$, là tập hợp tất cả các kết quả có thể xảy ra của phép thử. Mỗi phần tử của $\Omega$ là một kết quả có thể. Ta chỉ xét các phép thử có $\Omega$ hữu hạn; số phần tử của $\Omega$ kí hiệu là $n(\Omega)$.
\item \textbf{Biến cố} là một tập con của không gian mẫu: $A\subset\Omega$. Biến cố thường được mô tả bằng một mệnh đề hoặc liệt kê các phần tử.\\
Mỗi phần tử của $A$ là một \emph{kết quả thuận lợi} cho $A$; số kết quả thuận lợi là $n(A)$.
\item Biến cố xảy ra khi kết quả của phép thử thuộc tập $A$.
\item $\Omega$ là \textbf{biến cố chắc chắn} (luôn xảy ra); $\varnothing$ là \textbf{biến cố không thể} (không bao giờ xảy ra).
\item \textbf{Biến cố đối} của $A$, kí hiệu $\overline{A}$, là biến cố "$A$ không xảy ra":
\[ \boxed{\overline{A}=\Omega\setminus A},\qquad \boxed{n\left(\overline{A}\right)=n(\Omega)-n(A)}. \]
\item Công cụ đếm: quy tắc cộng, quy tắc nhân, sơ đồ cây, bảng; hoán vị $P_n=n!$, chỉnh hợp $A_n^k$, tổ hợp $C_n^k$.
\end{itemize}

\textbf{2. Phương pháp giải các dạng}

%% DANG 0D0?1-1
\textbf{Dạng 1. Các câu hỏi lý thuyết tổng hợp}
\begin{enumerate}
\item Nắm chắc các khái niệm phép thử, không gian mẫu, biến cố (là tập con của $\Omega$), biến cố chắc chắn, biến cố không thể và biến cố đối.
\item Xét mỗi khẳng định: đối chiếu với định nghĩa, hoặc lấy một phép thử cụ thể (gieo xúc xắc, tung đồng xu) để kiểm tra.
\end{enumerate}
Lưu ý: $\overline{A}$ là "$A$ không xảy ra", không phải biến cố "ngược nghĩa" theo cảm nhận. Ví dụ, khi gieo xúc xắc, biến cố đối của "ra số chấm lớn hơn $4$" là "ra số chấm không lớn hơn $4$", không phải "ra số chấm nhỏ hơn $4$".

%% DANG 0D0?1-2
\textbf{Dạng 2. Mô tả không gian mẫu, biến cố}
\begin{enumerate}
\item Xác định rõ phép thử: thực hiện mấy lần, có phân biệt thứ tự hay không, có hoàn lại hay không.
\item Viết $\Omega$ bằng cách liệt kê, dùng sơ đồ cây (phép thử nhiều giai đoạn) hoặc lập bảng (gieo $2$ xúc xắc).
\item Chọn trong $\Omega$ những kết quả thỏa mãn điều kiện để viết biến cố $A$; biến cố đối $\overline{A}$ gồm các kết quả còn lại.
\end{enumerate}
Lưu ý: tung $2$ đồng xu thì $SN$ và $NS$ là hai kết quả khác nhau. Gieo $2$ xúc xắc thì kết quả được ghi thành cặp có thứ tự $(i;j)$.

%% DANG 0D0?1-3
\textbf{Dạng 3. Đếm phần tử không gian mẫu, biến cố}
\begin{enumerate}
\item Đếm $n(\Omega)$ bằng quy tắc nhân, hoán vị, chỉnh hợp hoặc tổ hợp. Ví dụ: tung đồng xu $k$ lần thì $n(\Omega)=2^k$; gieo xúc xắc $k$ lần thì $n(\Omega)=6^k$; chọn $k$ phần tử từ $n$ phần tử (không thứ tự) thì $n(\Omega)=C_n^k$.
\item Đếm $n(A)$ trực tiếp bằng cách liệt kê hoặc dùng quy tắc cộng, quy tắc nhân.
\item Nếu $A$ có nhiều trường hợp hoặc có chữ "ít nhất", hãy đếm $n\left(\overline{A}\right)$ rồi tính $n(A)=n(\Omega)-n\left(\overline{A}\right)$.
\end{enumerate}
Lưu ý: dùng $A_n^k$ khi có thứ tự (xếp chỗ, phân công chức vụ, lập số) và $C_n^k$ khi không có thứ tự (chọn nhóm). Lấy có hoàn lại thì dùng quy tắc nhân $n^k$. Khi đếm theo từng trường hợp, kiểm tra để không đếm trùng.
